%% ch07 --- Liczba na chaos: wykładnik Lapunowa
%%
%% Napisany we wrześniu 2026. Każda liczba, której czytelnik nie policzy w
%% pamięci, pochodzi z code/measure/ch07.py, który importuje TE SAME funkcje,
%% które drukują listingi (code/ch07/), więc strona i listingi nie mogą się
%% rozejść. Orbity używają wyłącznie + - * /; logarytmy są uśredniane po stu
%% tysiącach składników i zapisywane z dwoma miejscami po przecinku, czego
%% żadna różnica ostatniego bitu między maszynami nie ruszy.
%%
%% CZEGO TEN ROZDZIAŁ NIE MOŻE POWIEDZIEĆ. Obserwacja, że dwa przebiegi się
%% rozchodzą, a ich odstęp mniej więcej się podwaja, należy do rozdziału 6;
%% ten rozdział zamienia ją w jedną liczbę i horyzont. Pełny diagram
%% bifurkacyjny, miejsce, gdzie zaczyna się chaos, zagęszczanie przejść i
%% największe okno należą do rozdziału 8. Wykładnik układu Lorenza należy do
%% rozdziału 9, czasy Lapunowa planet do rozdziału 13, horyzonty pogody do
%% rozdziału 14, a zaokrąglanie jako błąd początkowy do rozdziału 16.
%%
%% Bliźniak: chapters/en/ch07-lyapunov.tex. Podrozdział za podrozdziałem,
%% makro za makrem, liczba za liczbą; tools/parity.py je porównuje.

\chapter{Liczba na chaos: wykładnik Lapunowa}
\label{ch:lyapunov}
\index{wykładnik Lapunowa}\index{horyzont przewidywania}

\noindent
Rozdział~\ref{ch:butterfly} policzył dwie rzeczy o odwzorowaniu logistycznym
przy $r = 4$: odstęp między dwoma bliskimi przebiegami mniej więcej podwaja
się w każdym kroku, a każde poprawne miejsce po przecinku stanu początkowego
daje tylko kilka kroków prognozy więcej. Oba fakty znaleziono przez liczenie,
dla jednego odwzorowania. Pytanie, które za nimi stoi, zadasz każdemu
układowi w tej książce: jak szybko rosną małe błędy i jak daleko naprzód
pozwala to widzieć?

Ten rozdział zamienia liczenie w jedną liczbę, wykładnik Lapunowa, którą da
się zmierzyć dla każdego układu, a tę liczbę w horyzont: liczbę kroków, na
które przewidywanie jest dobre. Podczas
lektury trzymaj się tego pytania: gdybyś zmierzył stan początkowy układu
chaotycznego tysiąc razy dokładniej, o ile dalej mógłbyś go przewidywać?

\begin{outcomes}
\outcome{ustalić, jak bardzo jeden krok reguły rozciąga mały błąd, i sprawdzić
  to, uruchamiając dwie kopie programu}
\outcome{używać logarytmu naturalnego, żeby zamienić łańcuch mnożeń w sumę i
  z powrotem}
\outcome{obliczyć wykładnik Lapunowa odwzorowania logistycznego z orbity i
  odczytać jego znak}
\outcome{porównać obliczony wykładnik z wartością znaną dokładnie}
\outcome{zamienić wykładnik, błąd pomiaru i tolerancję w horyzont
  przewidywania}
\outcome{powiedzieć, co daje lepszy pomiar, licząc w krokach}
\end{outcomes}

\section{Krok po kroku}
\label{sec:ch07-stretch}
\index{rozciągnięcie}

\noindent
Rozdział~\ref{ch:butterfly} zauważył, że odstęp między dwoma przebiegami
rośnie nierówno: niektóre kroki prawie go czterokrotnie zwiększają, kilka go
zmniejsza, a podwaja się dopiero średnio. Zacznij od tej nierówności.
Uruchom odwzorowanie logistyczne przy $r = 4$ od $x = \num{0.3}$ i od
wartości o $\num{1e-10}$ wyższej i patrz, co każdy pojedynczy krok robi z
odstępem: podziel odstęp po kroku przez odstęp przed nim. Ten iloraz nazwijmy
\emph{rozciągnięciem} kroku.

\pyregion{code/ch07/stretch.py}{rule}{Odwzorowanie logistyczne i jego rozciągnięcie}{lst:ch07-rule}

\noindent
Funkcja \code{step} to odwzorowanie logistyczne z
rozdziału~\ref{ch:logistic}. Funkcja \code{stretch} to jego nachylenie,
$r(1 - 2x)$. Rozdział~\ref{ch:feedback} mierzył nachylenie, lekko przesuwając
$x$ i patrząc, jak daleko przesunie się wynik; dla tej reguły da się je
zapisać dokładnie, a pierwsze laboratorium sprawdza ten wzór przesuwaniem.
Reszta programu mierzy rozciągnięcie pierwszych ośmiu kroków na podstawie
dwóch przebiegów, obok nachylenia w miejscu, w którym jest pierwszy przebieg:

\pyregion{code/ch07/stretch.py}{table}{Pomiar rozciągnięcia ośmiu kroków}{lst:ch07-table}

\transcript{ch07-stretch}

\noindent
Dwie prawe kolumny się zgadzają. W jednym kroku malutki odstęp jest mnożony
przez nachylenie reguły w miejscu, w którym jesteś, więc rozciągnięcie można
odczytać z reguły, zamiast mierzyć je drugim przebiegiem. Ujemne nachylenie
dodatkowo przerzuca odstęp na drugą stronę, co nie zmienia jego wielkości;
tabela podaje wielkości.

A rozciągnięcia wcale nie przypominają stałej $2$. Krok $3$ zmniejsza odstęp
do mniej niż jednej trzeciej, a kroki $4$ i $5$ prawie go czterokrotnie
zwiększają. W kroku $3$ przebieg jest blisko $\num{0.5}$, gdzie parabola jest
płaska na szczycie; w krokach $4$ i $5$ jest blisko $1$ i blisko $0$, gdzie
parabola jest stroma. Chaos nie rozciąga równomiernie: w jednym miejscu
mocno, w innym łagodnie, a czasem ściska.

\begin{think}
Załóżmy, że jeden krok mnoży odstęp przez $4$, a następny przez $\tfrac14$.
Co stało się z odstępem w ciągu tych dwóch kroków? I jaki jeden czynnik,
użyty w obu krokach, dałby to samo?
\end{think}
\reveal{Nic się nie stało: $4 \times \tfrac14 = 1$, i odstęp wrócił tam,
gdzie był. Tym jednym czynnikiem jest $1$, a nie $\num{2.125}$, które
dostajesz, uśredniając $4$ i $\tfrac14$.}
Wzrost się składa: każdy krok mnoży to, co zostawił poprzedni. Typowe
rozciągnięcie nie jest więc zwykłą średnią rozciągnięć. To czynnik, który
użyty w każdym kroku daje ten sam wynik łączny.

\begin{trapbox}
\textbf{\enquote{Żeby poznać typowe rozciągnięcie na krok, uśrednij
rozciągnięcia.}} Sprawdź to na tabeli. Średnia ośmiu rozciągnięć wynosi
\val{ch07.eight.mean}, a odstęp pomnożony przez nią osiem razy urósłby mniej
więcej \val{ch07.eight.meanpow}-krotnie. Pomnóż zamiast tego osiem prawdziwych
rozciągnięć, a okaże się, że odstęp urósł \val{ch07.eight.product}-krotnie.
Kiedy coś się mnoży, dodawanie odpowiada na pytanie, którego nikt nie zadał;
średnia, której szukasz, też musi mnożyć.
\end{trapbox}

\section{Logarytm zamienia mnożenie w dodawanie}
\label{sec:ch07-log}
\index{logarytm!naturalny}

\noindent
Rozdział~\ref{ch:iteration} użył logarytmu, żeby odpowiedzieć na pytanie
\enquote{ile kroków, zanim\ldots?}. Ten rozdział potrzebuje jednego
konkretnego logarytmu, logarytmu \emph{naturalnego}, zapisywanego $\ln$, i
jednej rzeczy, którą on robi. Oto kilka jego wartości:

\begin{center}
\begin{tabular}{lcccccc}
\toprule
liczba & $\tfrac12$ & $1$ & $2$ & $4$ & $10$ & $e$ \\
\midrule
jej logarytm naturalny & $\val{ch07.ln.half}$ & $0$ & $\val{ch07.ln.two}$ &
  $\val{ch07.ln.four}$ & $\val{ch07.ln.ten}$ & $1$ \\
\bottomrule
\end{tabular}
\end{center}

\noindent
Spójrz na $2$ i $4$. Ponieważ $4 = 2 \times 2$, logarytm czwórki jest dwa
razy większy od logarytmu $2$. To wszystko, czego ten rozdział potrzebuje:
\[ \ln(a \times b) = \ln a + \ln b. \]
Słowami: logarytm iloczynu jest sumą logarytmów, więc mnożenie dwóch liczb
zamienia się w dodawanie ich logarytmów. Liczby większe od $1$ mają logarytmy
dodatnie, a liczby między $0$ a $1$ ujemne: $\tfrac12 \times 2 = 1$, więc
$\ln\tfrac12$ i $\ln 2$ sumują się do $\ln 1$, czyli do $0$. Ostatnia
kolumna zawiera $e$, mniej więcej \val{ch07.e}, liczbę, której logarytm
naturalny wynosi dokładnie $1$; odgrywa ona rolę, jaką $10$ odgrywa dla
logarytmu dziesiętnego. W Pythonie logarytm naturalny to \code{math.log}, a
\code{math.log(2)} zaczyna się od \val{ch07.ln.two.long}.

\begin{think}
Korzystając tylko z $\ln 2 \approx \val{ch07.ln.two}$, oblicz $\ln 8$. A ile
wynosi $\ln\tfrac14$?
\end{think}
\reveal{$8 = 2 \times 2 \times 2$, więc $\ln 8 = \ln 2 + \ln 2 + \ln 2$,
czyli mniej więcej \val{ch07.ln.eight}. A $\tfrac14 \times 4 = 1$, więc
$\ln\tfrac14$ wynosi $-\val{ch07.ln.four}$.}
Jeśli napisałeś $\num{2.07}$, zrobiłeś to dobrze: zaokrąglenie $\ln 2$
policzyłeś trzy razy.

Tak logarytm naprawia średnią. Jeśli odstęp jest w kroku mnożony przez $g$,
jego logarytm rośnie o $\ln g$: \emph{czynnik} na krok staje się
\emph{tempem} na krok, a tempa się dodają, więc można je uśredniać
zwyczajnie. Dla dwóch kroków z pytania tempa $\ln 4$ i $\ln\tfrac14$
sumują się do $0$, średnie tempo wynosi $0$, a czynnik, którego logarytm to
$0$, wynosi $1$: właściwa odpowiedź, znaleziona dodawaniem.

Żeby wrócić od tempa do czynnika, potrzebujesz liczby, której logarytmem jest
to tempo; zapisuje się ją $e^{\lambda}$, \enquote{$e$ do potęgi $\lambda$};
$\lambda$ to grecka litera lambda i oznacza tempo. Zatem $e^{0} = 1$,
$e^{1} = e$ i $e^{\ln 2} = 2$. Teraz osiem kroków z tabeli. Logarytmy ich
rozciągnięć mają średnią \val{ch07.eight.rate}, a
$e^{\val{ch07.eight.rate}}$ to mniej więcej \val{ch07.eight.geo}: średnio
odstęp rósł o czynnik około $2$ na krok, a \val{ch07.eight.geo} pomnożone
przez siebie osiem razy daje \val{ch07.eight.product}, które wyszło w tabeli.
\enquote{Mniej więcej się podwaja} stało się pomiarem.

\section{Wykładnik Lapunowa}
\label{sec:ch07-exponent}
\index{wykładnik Lapunowa!odwzorowania logistycznego}

\noindent
Osiem kroków to próbka. Prześledź zamiast tego orbitę przez sto tysięcy
kroków, weź logarytm naturalny wielkości rozciągnięcia w każdym punkcie,
który odwiedza, i uśrednij. Wynik to \emph{wykładnik Lapunowa}, nazwany na
cześć rosyjskiego matematyka Aleksandra Lapunowa, którego praca o
stabilności ruchu ukazała się w 1892 roku. Jeśli $s_1, s_2, \dots, s_N$ to
rozciągnięcia wzdłuż orbity,
\[ \lambda = \frac{\ln\abs{s_1} + \ln\abs{s_2} + \dots + \ln\abs{s_N}}{N}. \]
Słowami: $\lambda$ to średnia logarytmów rozciągnięć. Kreski oznaczają
\enquote{wielkość}, bo przerzucenie na drugą stronę nie czyni odstępu ani
większym, ani mniejszym.

Czytaj $\lambda$ jako tempo. W ciągu $N$ kroków mały błąd jest mnożony przez
mniej więcej $e^{\lambda N}$, więc znak mówi, z jakim układem masz do
czynienia:

\begin{itemize}
\item \textbf{$\lambda$ powyżej zera}: błędy rosną wykładniczo, a układ jest
  chaotyczny;
\item \textbf{$\lambda$ poniżej zera}: błędy maleją, a układ zapomina o
  swoich zaburzeniach;
\item \textbf{$\lambda$ równe zeru}: granica, na której błędy średnio ani nie
  rosną, ani nie maleją.
\end{itemize}

\begin{think}
Przy $r = 4$ odstęp mniej więcej podwaja się w każdym kroku. Jakiej wartości
$\lambda$ się spodziewasz?
\end{think}
\reveal{Mniej więcej $\ln 2$, czyli \val{ch07.ln.two}: odstęp mnożony w
każdym kroku przez $2$ rośnie w tempie $\ln 2$ na krok.}

\begin{think}
Przy $r = \num{2.8}$ orbita zatrzymuje się w punkcie stałym $1 - 1/r$ z
rozdziału~\ref{ch:logistic}. Tam nachylenie $r(1 - 2x)$ wynosi
$2 - r = -\num{0.8}$, a $\ln\num{0.8}$ to mniej więcej
\val{ch07.ln.pointeight}. Ile wynosi $\lambda$?
\end{think}
\reveal{Mniej więcej \val{ch07.ln.pointeight}. Kiedy orbita siedzi w punkcie
stałym, każdy składnik średniej to $\ln\num{0.8}$, więc średnia też. Jest
ujemna: zaburzenie maleje w każdym kroku o jedną piątą.}

\pyregion{code/ch07/exponent.py}{exponent}{Wykładnik Lapunowa przez uśrednianie}{lst:ch07-exponent}

\noindent
Pierwsza pętla wykonuje tysiąc kroków odwzorowania i je wyrzuca, żeby orbita
się ustaliła, zanim cokolwiek zostanie uśrednione. Druga dodaje logarytm
wielkości rozciągnięcia w każdym punkcie, a potem robi krok; ostatni wiersz
dzieli. To jest wzór powyżej i nic więcej. Uruchom go przy czterech
ustawieniach pokrętła $r$:

\pyregion{code/ch07/exponent.py}{run}{Wykładnik przy czterech ustawieniach r}{lst:ch07-run}

\transcript{ch07-exponent}

\noindent
Oba przewidywania się sprawdzają. Przy $r = \num{2.8}$ program znajduje
\val{ch07.lam.twoeight}. Przy $r = 4$ znajduje \val{ch07.lam.four}, a
dokładna wartość to $\ln 2 = \val{ch07.ln.two.long}$: średnia stu tysięcy
bardzo różnych logarytmów zgadza się z nią z dokładnością do $\num{0.001}$.
Przy $r = \num{3.2}$, gdzie orbita przeskakuje między dwiema wartościami, i
przy $r = \num{3.5}$, gdzie krąży między czterema, wykładnik też jest
ujemny: \val{ch07.lam.threetwo} i \val{ch07.lam.threefive}. Cykl, tak jak
punkt stały, zapomina o swoich zaburzeniach.

\plotfig{ch07-running-average}{Średnia bieżąca logarytmu rozciągnięcia
wzdłuż sześciu orbit: trzech startów przy $r = 4$ i trzech przy
$r = \num{3.7}$. Każda waha się przez kilkaset kroków, a potem się ustala, a
trzy starty przy każdym ustawieniu ustalają się na tej samej
liczbie.}{średnia, która się ustala}

\begin{rigourbox}
Dwie rzeczy są tu używane, ale nie dowodzone. Pierwsza: że średnia ustala się
na tej samej liczbie przy prawie każdym starcie. To twierdzenie, a
\enquote{prawie każdym} ma znaczenie: start dokładnie w $0$ zostaje w $0$ na
zawsze i daje średnią $\ln 4$. Przebieg na komputerze może to twierdzenie
zilustrować, jak robi to rysunek, ale nie może go dowieść. Druga: dlaczego
$r = 4$ daje dokładnie $\ln 2$. Zapisz $x$ jako $\sin^{2}\theta$ dla pewnego
kąta $\theta$, a jedna linijka szkolnej trygonometrii zamieni $4x(1 - x)$ w
$\sin^{2}(2\theta)$: przy $r = 4$ odwzorowanie to przebrane podwajanie kąta,
czyli odwzorowanie podwajające z rozdziału~\ref{ch:butterfly}, z kątem
liczonym w półobrotach. Podwajanie rozciąga każdy błąd kąta dokładnie $2$
razy, a zmiana przebrania nie zmienia średniego tempa.
\end{rigourbox}

\begin{lab}{l07_01_exponent}{Wykładnik z orbity}
Napisz \code{stretch}, czyli nachylenie $r(1 - 2x)$, oraz \code{exponent},
która pozwala orbicie się ustalić, a potem uśrednia wzdłuż niej logarytm
wielkości rozciągnięcia. Testy sprawdzają nachylenie przesuwaniem, tak jak
robił to rozdział~\ref{ch:feedback}, i porównują twój wykładnik z trzema
wartościami znanymi dokładnie: $\ln 2$ przy $r = 4$, $\ln\num{0.8}$ przy
$r = \num{2.8}$ i cyklem dwuokresowym przy $r = \num{3.2}$.
\end{lab}

\section{Jak daleko naprzód: horyzont przewidywania}
\label{sec:ch07-horizon}
\index{horyzont przewidywania}\index{czas Lapunowa}

\noindent
Po $n$ krokach mały błąd o wielkości $\delta$ (grecka litera delta) urasta do
mniej więcej $\delta\, e^{\lambda n}$. Przestajesz ufać przewidywaniu, gdy
jego błąd osiąga tolerancję $T$, największy błąd, z którym możesz żyć.
Horyzont to więc liczba kroków, przy której $\delta\, e^{\lambda n} = T$.
Zlogarytmuj obie strony: iloczyn staje się sumą, a logarytm $e^{\lambda n}$
to $\lambda n$, bo $e$ ma logarytm $1$ i jest mnożone $\lambda n$ razy.
Zostaje $\ln\delta + \lambda n = \ln T$, a stąd
\[ n \approx \frac{1}{\lambda}\,\ln\frac{T}{\delta}. \]
Słowami: horyzont to logarytm tego, ile razy błąd może urosnąć, podzielony
przez tempo, w jakim rośnie.

Liczba $1/\lambda$ nazywa się \emph{czasem Lapunowa}: to liczba kroków, w
ciągu których błędy rosną $e$ razy. Przy $r = 4$ wynosi $1/\ln 2$, mniej
więcej \val{ch07.lyap.time} kroku, a wzór mówi, że horyzont to czas Lapunowa
pomnożony przez $\ln(T/\delta)$.

\mermaidfig{ch07-three-steps}{Od reguły do horyzontu. Nachylenie daje
rozciągnięcie w każdym kroku, średnia jego logarytmu daje wykładnik, a
wykładnik razem z błędem i tolerancją daje horyzont.}{rozciągnięcie,
wykładnik, horyzont}

\begin{think}
Załóżmy, że model prognozy ma czas Lapunowa równy dwóm dniom, a jego
prognozy są dobre przez sześć. Nowa sieć przyrządów zmniejsza błąd jego
stanu początkowego dziesięciokrotnie. Na jak długo prognozy są dobre teraz:
na mniej więcej sześćdziesiąt dni, na mniej więcej dwanaście, czy krócej?
\end{think}
\reveal{Mniej więcej $\num{10.6}$ dnia. Dziesięciokrotnie mniejszy błąd
dodaje $\ln 10$ do $\ln(T/\delta)$, więc dodaje $\ln 10$ czasów Lapunowa:
$2 \times \val{ch07.ln.ten}$, czyli mniej więcej $\num{4.6}$ dnia.}

\begin{trapbox}
\textbf{\enquote{Dziesięć razy lepsze dane dają dziesięć razy dłuższe
prognozy.}} Dają stałą liczbę dodatkowych kroków, $\ln 10/\lambda$, bez
względu na to, jaki był wcześniej błąd. Przy $r = 4$ to \val{ch07.per.ten},
liczba, którą rozdział~\ref{ch:butterfly} policzył dla każdego miejsca po
przecinku, teraz wyprowadzona i prawdziwa dla każdego układu chaotycznego z
jego własnym $\lambda$. Horyzont rośnie z logarytmem twojej dokładności, a
nie z samą dokładnością: dwudzieste poprawne miejsce po przecinku daje nie
więcej niż pierwsze.
\end{trapbox}

\pyregion{code/ch07/horizon.py}{formula}{Wzór na horyzont}{lst:ch07-formula}

\noindent
Program zestawia ten wiersz z liczeniem, które wykonał
rozdział~\ref{ch:butterfly}: funkcja \code{counts} śledzi z tysiąca startów
dwa przebiegi oddalone o $\delta$, aż różnią się o więcej niż $T$. W
wydruku \code{1e-03} to pythonowy zapis $10^{-3}$, czyli jednej
tysięcznej.

\pyregion{code/ch07/horizon.py}{table}{Wzór kontra liczenie, dla czterech błędów stanu początkowego}{lst:ch07-count}

\transcript{ch07-horizon}

\noindent
Wzór i średnia z liczenia zgadzają się z dokładnością do kroku; liczenie
wychodzi odrobinę później, bo odstęp sprawdza się tylko po pełnych krokach,
a gdy przestaje być mały, rośnie wolniej. Dwie ostatnie kolumny są
ostrzeżeniem: poszczególne starty się różnią, niektóre bardzo, bo horyzont,
tak jak wykładnik, jest średnią. Ale spójrz w dół wierszy. Każda tysiąckrotna
poprawa dodaje te same dziesięć kroków; w pomiarze wiersze dzieli średnio
\val{ch07.hz.gained} kroku, wobec $\ln 1000/\ln 2$ ze wzoru, czyli
\val{ch07.ln.thousand} podzielone przez \val{ch07.ln.two}:
\val{ch07.extra.thousand}.

\plotfig{ch07-horizon-line}{Horyzont w zależności od błędu stanu
początkowego, w skali logarytmicznej, z lepszymi pomiarami po prawej.
Liczenie (kropki i pasmo, w którym mieści się większość startów) idzie
wzdłuż prostej ze wzoru: każdy czynnik dziesięć w dokładności daje tę samą
liczbę dodatkowych kroków.}{horyzont a błąd stanu początkowego}

\begin{think}
Komputer przechowuje liczbę z dokładnością do mniej więcej szesnastu cyfr
znaczących. Załóżmy, że potrafisz tak dobrze zmierzyć stan początkowy
odwzorowania przy $r = 4$, więc $\delta = 10^{-16}$. Przy $T = \num{0.1}$ jak
daleko naprzód mógłbyś mniej więcej przewidywać? I ilu miejsc po przecinku
stanu początkowego wymagałoby tysiąc kroków?
\end{think}
\reveal{Mniej więcej \val{ch07.hz.sixteen} kroków: piętnaście czynników
dziesięć, po \val{ch07.per.ten} kroku każdy. Tysiąc kroków wymagałoby stanu
początkowego z dokładnością do \val{ch07.digits.thousand} miejsc po
przecinku; nawet sto kroków wymaga \val{ch07.digits.hundred}.}

\begin{trapbox}
\textbf{\enquote{Większy komputer w końcu uczyni każdy układ
przewidywalnym.}} Szybszy komputer nie zna stanu początkowego ani trochę
lepiej, a horyzont zależy od stanu początkowego tylko przez logarytm jego
błędu. Żeby podwoić horyzont, musisz podwoić liczbę poprawnych cyfr, a
najdokładniejsze pomiary, jakich kiedykolwiek dokonano, za pomocą zegarów
atomowych, sięgają mniej niż dwudziestu. Komputer w dodatku zaokrągla w
każdym kroku, co jest nowym błędem, jego własnym;
rozdział~\ref{ch:computers} mierzy, ile to kosztuje.
\end{trapbox}

\noindent
Dlatego prognozy pogody poprawiają się powoli. Satelity, gęstsze obserwacje
i szybsze komputery wielokrotnie zmniejszyły błąd zmierzonego stanu
atmosfery, a każdy taki czynnik kupił tę samą stałą porcję wyprzedzenia;
zwykle przytacza się poprawę o mniej więcej jeden dzień użytecznej prognozy
na dekadę. Atmosfera nie jest odwzorowaniem logistycznym, a
rozdział~\ref{ch:weather} mówi, co jest w niej inne, ale arytmetyka
logarytmu stosuje się do niej tak samo.

\begin{worldbox}
Czasem Lapunowa astronomowie opisują, jak daleko naprzód da się śledzić
orbitę. W 1988 roku Gerald Sussman i Jack Wisdom prześledzili Plutona na
komputerze zbudowanym specjalnie w tym celu, stwierdzili, że jego ruch jest
chaotyczny, i oszacowali jego czas Lapunowa na mniej więcej dwadzieścia
milionów lat; historia samych planet należy do
rozdziału~\ref{ch:mechanics}. Sama ta liczba nie jest horyzontem. Horyzont
to czas Lapunowa pomnożony przez $\ln(T/\delta)$, a ten logarytm wynosi
tylko mniej więcej \val{ch07.ln.tenbillion} nawet wtedy, gdy błąd stanu
początkowego jest dziesięć miliardów razy mniejszy od tolerancji, więc
horyzont to co najwyżej kilkadziesiąt czasów Lapunowa.
\end{worldbox}

\begin{lab}{l07_02_horizon}{Jak daleko naprzód i jakim kosztem}
Napisz \code{horizon}, czyli wzór z tego podrozdziału; \code{extra\_steps},
czyli liczbę kroków zyskanych dzięki zmierzeniu stanu początkowego zadaną
liczbę razy dokładniej; oraz \code{digits\_needed}, czyli liczbę miejsc po
przecinku, jakiej wymaga prognoza zadanej długości. Testy sprawdzają, że
tysiąckrotna poprawa daje tę samą liczbę kroków bez względu na to, od jakiego
błędu zaczynasz.
\end{lab}

\section{Chaos przychodzi i odchodzi}
\label{sec:ch07-sweep}
\index{wykładnik Lapunowa!w zależności od r}

\noindent
Mając jedną liczbę na każde ustawienie pokrętła, możesz kręcić pokrętłem i
ją obserwować. Rozdział~\ref{ch:logistic} znalazł punkt stały przy
$r = \num{2.8}$ i cykl dwuokresowy przy $r = \num{3.2}$; między nimi, przy
$r = 3$, punkt stały ustępuje.

\begin{think}
W punkcie stałym nachylenie wynosi $2 - r$. Ile wynosi przy $r = 3$ i co to
oznacza dla $\lambda$?
\end{think}
\reveal{$2 - 3 = -1$, a $\ln 1 = 0$, więc $\lambda = 0$: zaburzenie ani się
nie rozciąga, ani nie kurczy, tylko przerzuca. Uruchomiony przy $r = 3$
program znajduje wartość odległą od zera o mniej niż $\num{0.001}$.}
Zero to granica i wskazuje miejsca, gdzie jeden rodzaj zachowania przechodzi
w inny. Żeby takie miejsca znaleźć, policz wykładnik przy wielu ustawieniach
$r$. Listing używa funkcji \api{lyapunov} z biblioteki, która jest pętlą z
\code{exponent}, z regułą i jej rozciągnięciem przekazanymi jako argumenty,
więc działa dla każdego odwzorowania:

\pyregion{code/ch07/sweep.py}{sweep}{Wykładnik przy szesnastu ustawieniach r}{lst:ch07-sweep}

\transcript{ch07-sweep}

\noindent
Czytaj środkową kolumnę od góry. Jest ujemna aż do $r = \num{3.56}$ i
dodatnia od $r = \num{3.60}$, więc chaos pojawia się gdzieś pomiędzy. Potem
rośnie ku $\ln 2$ przy $r = 4$, z wyjątkiem $r = \num{3.84}$, gdzie znów jest
ujemna: to wąski zakres pokrętła wewnątrz chaosu, w którym orbita znalazła
cykl i znów zapomina o zaburzeniach. Tak rzadka tabela przeskakuje większość
takich zakresów; rysunek nie.

\plotfig{ch07-exponent-r}{Wykładnik Lapunowa w zależności od pokrętła $r$.
Jest ujemny wszędzie tam, gdzie orbita się ustala, dotyka zera w każdym
miejscu przejścia, jest dodatni tam, gdzie odwzorowanie jest chaotyczne, i
znów spada poniżej zera w wąskich oknach. Czerwone kropki to cztery wartości
z listingu wykładnika.}{wykładnik w zależności od $r$}

\noindent
Przed początkiem chaosu krzywa nigdy nie wznosi się ponad zero. Dotyka zera
przy $r = 3$ i jeszcze raz trochę poniżej $r = \num{3.5}$, za każdym razem,
gdy cykl przechodzi w cykl dwa razy dłuższy. Dalej jest przeważnie dodatnia
i poprzebijana skierowanymi w dół szpilkami: oknami, w których wraca
porządek. Gdzie dokładnie zaczyna się chaos, dlaczego przejścia się
zagęszczają i które okno jest najszersze, to pytania
rozdziału~\ref{ch:bifurcations}. Odpowiedzią tego rozdziału jest pomiar,
który przy każdym ustawieniu odróżnia porządek od chaosu: jedna liczba i jej
znak.

\begin{summarybox}
\begin{itemize}
\item Jeden krok mnoży malutki błąd przez \textbf{rozciągnięcie}, czyli
  nachylenie reguły w miejscu, w którym jesteś. Wzdłuż chaotycznej orbity
  rozciągnięcia bardzo się różnią, a ich zwykła średnia jest złym
  podsumowaniem.
\item \textbf{Logarytm naturalny} zamienia mnożenie w dodawanie:
  $\ln(a \times b) = \ln a + \ln b$, $\ln 1 = 0$, $\ln 2 \approx
  \val{ch07.ln.two}$ i $\ln e = 1$. Czynnik wzrostu $g$ na krok staje się
  tempem $\ln g$, a tempa można uśredniać.
\item \textbf{Wykładnik Lapunowa} $\lambda$ to średnia z
  $\ln\abs{\text{rozciągnięcie}}$ wzdłuż orbity: dodatni dla chaosu, ujemny
  dla porządku, zero na granicy. Dla odwzorowania logistycznego przy $r = 4$
  wynosi $\ln 2$, i obliczenie się z tym zgadza.
\item \textbf{Horyzont przewidywania} to mniej więcej
  $\frac{1}{\lambda}\ln(T/\delta)$, a $1/\lambda$ to \textbf{czas Lapunowa}.
  Stan początkowy zmierzony tysiąc razy dokładniej daje $\ln 1000/\lambda$
  dodatkowych kroków, przy $r = 4$ mniej więcej dziesięć: stałą porcję, a nie
  tysiąc razy więcej.
\item W zależności od $r$ wykładnik pokazuje, jak \textbf{chaos przychodzi i
  odchodzi}: jest ujemny tam, gdzie orbita się ustala, zerowy w każdym
  miejscu przejścia, dodatni w chaosie i znów ujemny w wąskich oknach.
\end{itemize}
\end{summarybox}

\canyou

\begin{checkpoint}
\item Odstęp jest mnożony przez $3$, potem przez $\tfrac13$, potem przez $3$,
  potem przez $\tfrac13$. O ile w sumie urósł i ile wynosi zwykła średnia tych
  czterech czynników?
  \answerto{Wcale nie urósł: iloczyn wynosi $1$. Zwykła średnia czynników to
  $\tfrac53$, co sugeruje wzrost, którego nie było.}
\item Korzystając z $\ln 2 \approx \val{ch07.ln.two}$ i
  $\ln 10 \approx \val{ch07.ln.ten}$, oblicz $\ln 20$ i $\ln\num{0.2}$.
  \answerto{$\ln 20 = \ln 2 + \ln 10 \approx \num{2.99}$. A
  $\num{0.2} \times 10 = 2$, więc $\ln\num{0.2} = \ln 2 - \ln 10 \approx
  -\num{1.61}$.}
\item Orbita ustala się w punkcie stałym, w którym nachylenie wynosi
  $-\num{0.5}$. Ile wynosi wykładnik Lapunowa i co mówi jego znak?
  \answerto{$\ln\num{0.5} \approx -\num{0.69}$. Jest ujemny: błędy maleją, w
  każdym kroku o połowę.}
\item Układ ma $\lambda = \num{0.5}$ na krok. Błąd $\num{0.001}$ musi pozostać
  poniżej $\num{0.1}$. Na mniej więcej ile kroków przewidywanie jest dobre?
  \answerto{Mniej więcej $9$: $\ln(\num{0.1}/\num{0.001}) = \ln 100 = 2\ln 10
  \approx \num{4.6}$, podzielone przez $\num{0.5}$.}
\item W tym samym układzie mierzysz stan początkowy sto razy dokładniej. Ile
  kroków zyskujesz? A jeśli zrobisz to jeszcze raz?
  \answerto{Mniej więcej $9$ kroków za każdym razem:
  $\ln 100/\num{0.5} \approx \num{9.2}$. Każda stukrotna poprawa daje tę samą
  stałą liczbę kroków.}
\item Dlaczego szybszy komputer sam w sobie prawie nie wydłuży prognozy
  układu chaotycznego?
  \answerto{Bo horyzont zależy od błędu stanu początkowego tylko przez jego
  logarytm, a szybszy komputer nie mierzy stanu początkowego ani trochę
  lepiej. Podwojenie horyzontu wymaga dwa razy więcej poprawnych cyfr stanu
  początkowego.}
\item Przy $r = \num{3.84}$ wykładnik jest ujemny, choć przy $r = \num{3.80}$
  i $r = \num{3.88}$ jest dodatni. Co się tam dzieje?
  \answerto{Okno: wewnątrz chaotycznego zakresu orbita ustaliła się w cyklu,
  więc błędy znów maleją, a wykładnik spada poniżej zera.}
\end{checkpoint}
